% !TEX TS-program = pdflatex
% !TeX program = pdflatex
% !TEX encoding = UTF-8
% !TEX spellcheck = en_US

\documentclass[xcolor=table]{beamer}


%\usepackage{fullpage}
%\usepackage[left=2.8cm,right=2.2cm,top=2 cm,bottom=2 cm]{geometry}
\setbeamersize{text margin left=10pt,text margin right=10pt}



%\usepackage[T1]{fontenc}
\usepackage{textcomp}
\usepackage[utf8]{inputenc}
%\usepackage[french,english]{babel}
\usepackage{arabtex}
\usepackage{txfonts}
\usepackage{tipa}
\usepackage[]{graphicx}
\usepackage{multirow}
\usepackage{hyperref}
\usepackage{colortbl}
\usepackage{tabularray}
\usepackage{listingsutf8}
\usepackage{wrapfig}
\usepackage{multicol}
\usepackage[export]{adjustbox} %for images in table, also for frame
\usepackage[many]{tcolorbox}
\usepackage{wasysym}
\usepackage[lined]{algorithm2e}
\usepackage{alltt} %verbatim with commands
\usepackage{longtable}
\usepackage{tabu}

\usepackage{CJKutf8}

\usepackage{amsmath,amssymb, amsfonts} 

%\usepackage{newtxtext,newtxmath}


%\usepackage{acolor}


\definecolor{lightblue}{HTML}{D0D2FF}
\definecolor{lightyellow}{HTML}{FFFFAA}
\definecolor{darkblue}{HTML}{0000BB}
\definecolor{olivegreen}{HTML}{006600}
\definecolor{darkgreen}{HTML}{008B45} %009B55
\definecolor{violet}{HTML}{6600CC}
\definecolor{deeppink}{HTML}{FF1493}
\definecolor{orangey}{HTML}{FFBB00}


\makeatletter
\def\beamer@calltheme#1#2#3{%
	\def\beamer@themelist{#2}
	\@for\beamer@themename:=\beamer@themelist\do
	{\IfFileExists{\beamer@themelocation/#3\beamer@themename.sty}{%
	\usepackage[{#1}]{\beamer@themelocation/#3\beamer@themename}}{\usepackage[{#1}]{#3\beamer@themename}}
		
	}%
}

\def\usefolder#1{
	\def\beamer@themelocation{#1}
}
\def\beamer@themelocation{}
\makeatother

\usefolder{../../extra/beamer-themes}

\usetheme{Karim} % Antibes Boadilla Warsaw Karim

\beamertemplatenavigationsymbolsempty

\let\oldcite\cite
\renewcommand{\cite}[1]{{\bfseries\color{blue}\oldcite{#1}}}

\tcbuselibrary{listings}


%\renewcommand{\baselinestretch}{1.5}

\def\supit#1{\raisebox{0.8ex}{\small\it #1}\hspace{0.05em}}

\AtBeginSection{%
	\begin{frame}
		\sectionpage
	\end{frame}
}

\newcommand{\rottext}[2]{%
	\rotatebox{90}{%
	\begin{minipage}{#1}%
		\raggedleft#2%
	\end{minipage}%
	}%
}


\institute{ %
Laboratoire de la Communication dans les Systèmes Informatiques (LCSI)
\par
École  nationale Supérieure d'Informatique (ESI), Algiers, Algeria
}
\author[Abdelkrime Aries (2023/2024)] %
{Abdelkrime Aries}
%\titlegraphic{\includegraphics[height=1cm]{../img/esi-logo.png}%\hspace*{4.75cm}~


\date{Academic year: 2023/2024} %\today

\titlegraphic{%
	\includegraphics[height=1cm]{../../img/misc/esi-logo.png}%
	\hspace{2cm}%
	\includegraphics[height=1cm]{../../img/misc/lcsi-logo.png}%
	\hspace{2cm}%
	\includegraphics[height=1.5cm]{../../img/misc/esi.nlp-logo.png}%
}

%\setbeamertemplate{headline}{}

\newcommand{\kurl}[1]{{\scriptsize\bfseries\color{orangey}\url{#1}}}

\newcommand{\keyword}[1]{\textcolor{red}{\bfseries\itshape #1}}
\newcommand{\expword}[1]{\textcolor{olivegreen}{#1}}
\newcommand{\optword}[1]{\textcolor{violet}{\bfseries #1}}

\makeatletter
\newcommand\mysphere{%
	\parbox[t]{10pt}{\raisebox{0.2pt}{\beamer@usesphere{item projected}{bigsphere}}}}
\makeatother

%\let\oldtabular\tabular
%\let\endoldtabular\endtabular
%\renewenvironment{tabular}{\rowcolors{2}{white}{lightblue}\oldtabular\rowcolor{blue}}{\endoldtabular}


%\NoAutoSpacing %french autospacing after ":"

\def\graphpath{}

\newcommand{\changegraphpath}[1]{\def\graphpath{#1}}


\newcommand{\vgraphpage}[2][.84\textheight]{%
%	\begin{center}%
		\includegraphics[height=#1]{\graphpath #2}%
%	\end{center}%
}

\newcommand{\hgraphpage}[2][\textwidth]{%
%	\begin{center}%
		\includegraphics[width=#1]{\graphpath #2}%
%	\end{center}%
}

\newcommand{\graphpage}[2][]{%
	\includegraphics[#1]{\graphpath #2}%
}

\bibliographystyle{apalike}

\newcommand{\insertbibliography}[2]{
	\appendix
	\section*{Bibliography}
	\nocite{#2}
%	\makeatletter % to change template
%	\setbeamertemplate{headline}[default] % not mandatory, but I though it was better to set it blank
%	\def\beamer@entrycode{\vspace*{-\headheight}} % here is the part we are interested in :)
%	\makeatother
	\begin{multicols*}{2}[\frametitle{\insertsection} \usebeamertemplate*{frametitle}]%\usebeamertemplate*{frametitle}\frametitle{Références}
		\tiny
		\bibliography{#1}
	\end{multicols*}
}

\definecolor{my-grey}{RGB}{233, 233, 233}

\newcommand{\insertlicence}{
	\begin{frame}[plain]
	\frametitle{License: CC-BY 4.0}

	\begin{tcolorbox}[colback=cyan,
		colframe=cyan,  
		arc=0pt,outer arc=0pt,
		valign=top, 
		halign=center,
		width=\textwidth]
		
		\includegraphics[width=.5cm]{../../img/licence/cc_icon_white_x2.png}
		\includegraphics[width=.5cm]{../../img/licence/attribution_icon_white_x2.png}
		
		\color{white}
		\bfseries Attribution 4.0 International (CC BY 4.0) \\
		\tiny \url{https://creativecommons.org/licenses/by/4.0/deed.en}
		
	\end{tcolorbox}\vspace{-.5cm}
	\begin{tcolorbox}[colback=my-grey,
		colframe=my-grey,  
		center, arc=0pt,outer arc=0pt,
		valign=top, 
		halign=left,
		width=\textwidth]
		
		\tiny
		
		\begin{center}
			\bfseries\Large
			You are free to:
		\end{center}
		
		\begin{minipage}{0.83\textwidth}
			\begin{itemize}
				\item[] \textbf{Share} — copy and redistribute the material in any medium or format
				\item[] \textbf{Adapt} — remix, transform, and build upon the material
				for any purpose, even commercially.
			\end{itemize}
		
			\vspace{12pt}
		
			\textit{The licensor cannot revoke these freedoms as long as you follow the license terms.}
		\end{minipage}
		\begin{minipage}{0.15\textwidth}
			\includegraphics[width=\textwidth]{../../img/licence/FreeCulturalWorks_seal_x2.jpg}
		\end{minipage}
	
		
		\begin{center}
			\bfseries\Large
			Under the following terms:
		\end{center}
		
		\begin{itemize}
			\item[] \textbf{Attribution} — You must give appropriate credit, provide a link to the license, and indicate if changes were made. You may do so in any reasonable manner, but not in any way that suggests the licensor endorses you or your use.
			\item[] \textbf{No additional restrictions} — You may not apply legal terms or technological measures that legally restrict others from doing anything the license permits.
		\end{itemize}
		
	\end{tcolorbox}
	
%	\begin{center}
%		\bfseries Attribution 4.0 International (CC BY 4.0)
%		\url{https://creativecommons.org/licenses/by/4.0/deed.fr}
%	\end{center}

%	\tiny
%
%	Vous êtes autorisé à : 
%	\begin{itemize}
%		\item \textbf{Partager} — copier, distribuer et communiquer le matériel par tous moyens et sous tous formats
%		\item \textbf{Adapter} — remixer, transformer et créer à partir du matériel
%	\end{itemize}
%	
%	Selon les conditions suivantes : 
%	\begin{itemize}
%		\item \textbf{Attribution} — Vous devez créditer l'Œuvre, intégrer un lien vers la licence et indiquer si des modifications ont été effectuées à l'Oeuvre. Vous devez indiquer ces informations par tous les moyens raisonnables, sans toutefois suggérer que l'Offrant vous soutient ou soutient la façon dont vous avez utilisé son Oeuvre.
%		\item \textbf{Pas d'Utilisation Commerciale} — Vous n'êtes pas autorisé à faire un usage commercial de cette Oeuvre, tout ou partie du matériel la composant. 
%		\item \textbf{Pas de restrictions complémentaires} — Vous n'êtes pas autorisé à appliquer des conditions légales ou des mesures techniques qui restreindraient légalement autrui à utiliser l'Oeuvre dans les conditions décrites par la licence.
%	\end{itemize}

	\end{frame}
}

\settowidth{\leftmargini}{\usebeamertemplate{itemize item}}
\addtolength{\leftmargini}{\labelsep}

\AtBeginDocument{
	\newcolumntype{L}[2]{>{\vbox to #2\bgroup\vfill\flushleft}p{#1}<{\egroup}} 
	
	\begin{frame}[plain]
		\maketitle
	\end{frame}

	\insertlicence
}


% needs etoolbox; to break links after -
\appto\UrlBreaks{\do\-}


\makeatletter
\newcommand{\xRightarrow}[2][]{\ext@arrow 0359\Rightarrowfill@{#1}{#2}}
\makeatother


\usefonttheme{structurebold}
%\usefonttheme{professionalfonts}


% figure caption
\setlength\abovecaptionskip{2pt}

% equations
\AtBeginDocument{
\abovedisplayskip=0pt
\abovedisplayshortskip=0pt
\belowdisplayskip=0pt
\belowdisplayshortskip=0pt
}


\title[ESI - NLP\_master: 01- Introduction]%
{Natural Language Processing (master)\\Chapter 01\\Introduction to NLP} 

\changegraphpath{../../img/intro/}

\begin{document}
	
\begin{frame}
\frametitle{Natural Language Processing}
\framesubtitle{Introduction: Definitions}

\begin{itemize}
	\item \keyword{NLP}: Natural Language Processing.
	\item all methods that make human language accessible to computers.
\end{itemize}
\begin{minipage}{0.78\textwidth}
\begin{itemize}
	\item A subfield of:
	\begin{itemize}
		\item \optword{Linguistics}: scientific study of human language.
		\item \optword{Computer science}: study of computation, automation and information.
		\item \optword{Artificial intelligence}: the science and engineering of
		making intelligent machines (John McCarthy, 1955).
	\end{itemize}
\end{itemize}
\end{minipage}
\begin{minipage}{0.20\textwidth}
\hgraphpage[\textwidth]{NLP.pdf}
\end{minipage}
\end{frame}

\begin{frame}
\frametitle{Natural Language Processing}
\framesubtitle{Introduction: Motivation}

\begin{itemize}
	\item Large volume of textual data;
	
	\item Texts: unstructured data;
	
	\item Communications;
	
	\item Popularity among companies;
	
	\item ...
\end{itemize}

\end{frame}

\begin{frame}
\frametitle{Natural Language Processing}
\framesubtitle{Introduction: Some humor}

\begin{center}
	\vgraphpage{humor/humor.jpg}
\end{center}

\end{frame}

\begin{frame}
\frametitle{Natural Language Processing}
\framesubtitle{Introduction: Plan}

\begin{multicols}{2}
%	\small
\tableofcontents
\end{multicols}
\end{frame}

%===================================================================================
\section{History}
%===================================================================================

\begin{frame}
\frametitle{Introduction to NLP}
\framesubtitle{History}

\hgraphpage{history.pdf}

\end{frame}

\subsection{AI birth and golden era}

\begin{frame}
\frametitle{Introduction to NLP: History}
\framesubtitle{AI birth and golden era: The 1950s}

\begin{itemize}
	\item \optword{1951}: Shannon exploited probabilistic models of natural languages \cite{1951-shannon}.
	\item \optword{1954}: Georgetown-IBM experiment to automatically translate 60 sentences from Russian to English.
	\item \optword{1956}: Chomsky described a hierarchy of classes of formal grammars; known as \keyword{Chomsky hierarchy}.
	\item \optword{1958}: Luhn (IBM) experimented on automatic text summarization by extraction \cite{1958-luhn}.
\end{itemize}

\end{frame}

\begin{frame}
\frametitle{Introduction to NLP: History}
\framesubtitle{AI birth and golden era: The 1960s}

\begin{itemize}
	\item \optword{1961}: Development of the first automatic parser at UPenn. \cite{1961-joshi,1962-harris} 
	\item \optword{1964}: Weizenbaum has developed \keyword{ELIZA}, a simulation of a psychotherapist within the MIT AI lab.
	\item \optword{1964}: Bobrow has developed \keyword{STUDENT} which is designed for reading and solving word problems found in high school algebra books. \cite{1964-bobrow}.
	\item \optword{1967}: Brown corpus; the first electronic corpus.
\end{itemize}

\end{frame}

\subsection{AI Winter}

\begin{frame}
\frametitle{Introduction to NLP: History}
\framesubtitle{AI Winter: The 1970s}

\begin{itemize}
	\item \optword{1971}: Winograd (MIT) has developed \keyword{SHRDLU}, a natural language understanding program \cite{1971-winograd}.
	\item \optword{1972}: Colby (Stanford) has created \keyword{PARRY}, a chatbot that simulates a person with paranoid schizophrenia.
	\item \optword{1975}: \keyword{MARGIE} a system that makes inferences and paraphrases from sentences using language's conceptual representation.
	\item \optword{1975}: \keyword{DRAGON}, a system for automatic speech recognition using hidden Markov models \cite{1975-baker}.
\end{itemize}

\end{frame}

\begin{frame}
\frametitle{Introduction to NLP: History}
\framesubtitle{AI Winter: The 1980s}

\begin{itemize}
	\item \optword{1980}: \keyword{KL-One}, knowledge representation for processing syntax and semantics \cite{1980-bobrow}.
	\item \optword{1986}: \keyword{TRUMP}, language analyzer using a lexical database \cite{1986-jacobs}.
	\item \optword{1987}: \keyword{HPSG}; head-driven phrase structure grammar \cite{1987-sag-pollard}.
	\item \optword{1987}: \keyword{MUC}; Message Understanding Conference which is financed by \keyword{DARPA}.
	\item \optword{1988}: Using hidden Markov models for part of speech tagging \cite{1988-church}.
	\item Symbolic solutions on speech processing and natural language generation.
\end{itemize}

\end{frame}

\subsection{AI Spring}

\begin{frame}
\frametitle{Introduction to NLP: History}
\framesubtitle{AI Spring: The 1990s}

\begin{itemize}
	\item \optword{1990}: Statistical approach for machine translation \cite{1990-brown-al}.
	\item \optword{1993}: \keyword{Penn Treebank}, annotated English corpus \cite{1993-marcus-al}.
	\item \optword{1995}: \keyword{Wordnet}, lexical database of English \cite{1995-miller}.
	\item \optword{1996}: \keyword{SPATTER}, statistical lexical analyzer based on decision trees \cite{1996-magerman}.
	\item Popularity of statistical methods and empirical evaluation.
\end{itemize}

\end{frame}

\begin{frame}
\frametitle{Introduction to NLP: History}
\framesubtitle{AI Spring: The 2000s}

\begin{itemize}
	\item \optword{2003}: Neural language models \cite{2003-bengio-al}.
	\item \optword{2006}: \keyword{Watson} (IBM), question/answering system.
	\item Use of unsupervised and semi-supervised learning as alternatives to purely supervised learning.
	\item Shift focus to semantic tasks.
\end{itemize}

\end{frame}

\begin{frame}
\frametitle{Introduction to NLP: History}
\framesubtitle{AI Spring: The 2010s}

\begin{itemize}
	\item \optword{2011}: \keyword{Siri} (Apple),  a personal digital assistant (PDA). 
	It was followed by \keyword{Alexa} (Amazon, \optword{2014}) and \optword{Google Assistant} (\optword{2016}).
	\item \optword{2014}: Word embedding \cite{2014-lebret-collobert}.
	\item \optword{2018}: Emergence of contextual representations (pre-trained language models) such as \keyword{ULMfit} (fast.ai) \cite{2018-howard-ruder}, \keyword{ELMO} (AllenNLP) \cite{2018-peters-al}, \keyword{GPT} (OpenAI) \cite{2018-radford-al}, \keyword{BERT} (Google) \cite{2018-devlin-al}, \keyword{XLM} (Facebook) \cite{2019-lample-conneau}.
\end{itemize}

\end{frame}

\begin{frame}
	\frametitle{Introduction to NLP: History}
	\framesubtitle{AI Spring: The 2020s}
	
	\begin{itemize}
		\item \optword{2020}: Interpretability of neural NLP \cite{2020-belinkov-al,2020-kim-al,2020-jacovi-goldberg}.
		\item \optword{2021}: Ethics of NLP \cite{2021-garridoMunoz-al,2021-field-al}.
		\item \optword{2022}: Polyglot language models such as \keyword{Polyglot Prompt} \cite{2022-fu-al}.
		\item \optword{2022}: Edge NLP (do more with less) \cite{2022-guo-al}.
	\end{itemize}
	
\end{frame}

%===================================================================================
\section{Language processing levels}
%===================================================================================

\begin{frame}
\frametitle{Introduction to NLP}
\framesubtitle{Language processing levels}

\begin{center}
	\vgraphpage{levels.pdf}
\end{center}

\end{frame}

\subsection{Phonetics, phonology et orthography}

\begin{frame}
\frametitle{Introduction to NLP: Levels}
\framesubtitle{Phonetics, phonology, orthography: Phonetics}

\begin{itemize}
	\item Study of the sounds or phones produced by the human vocal apparatus.
	\item Branches:
	\begin{itemize}
		\item \optword{Articulatory phonetics} (most anatomical and physiological branch) describes how vowels and consonants are produced in various parts of the mouth and throat.
		\item \optword{Acoustic phonetics} (the branch that has the closest affinities with physics) studies the sound waves that transmit vowels and consonants through the air from speaker to listener.
		\item \optword{Auditory phonetics} (the branch that most interests psychologists) examines how the listener's brain decodes sound waves into vowels and consonants originally intended by the speaker.
	\end{itemize}
\end{itemize}

\end{frame}

\begin{frame}
\frametitle{Introduction to NLP: Levels}
\framesubtitle{Phonetics, phonology, orthography: Phonetics (Articulation points of consonants)}

\begin{minipage}{0.55\textwidth}
\begin{itemize}
	\item \optword{labial}: Lips. E.g., \expword{\textipa{[b], [p], [m], [f], [v]}}.
	\item \optword{apical}: Tongue's tip or its anterior part. 
	E.g., \expword{%
		\textipa{[t], [d], [n], [r], }
		\RL{^s} \textipa{[S],} 
		\RL{_t} \textipa{[T],} 
		\RL{_d} \textipa{[D]}
	}
	\item \optword{dorsal}: Posterior part of the tongue. E.g., \expword{\textipa{[c], [k], [g], [q],} \RL{.g} \textipa{[G]}}.
	\item \optword{pharyngeal}: Pharynx. 
	E.g., \expword{\RL{.h} \textipa{[\*h],} \RL{`} \textipa{[Q]}}.
	\item \optword{glottal}: Glottis. 
	E.g., \expword{\textipa{[h],} \RL{'} \textipa{[P]}}.
\end{itemize}
\end{minipage}
\begin{minipage}{0.43\textwidth}
	\begin{figure}
		\hgraphpage{vocaltract.pdf}\caption{Vocal tract \cite{2009-ball}}
	\end{figure}
\end{minipage}

\end{frame}

\begin{frame}
\frametitle{Introduction to NLP: Levels}
\framesubtitle{Phonetics, phonology, orthography: Phonetics (Articulation modes of consonants)}

\begin{itemize}
	\item \optword{occlusive}: complete blockage of airflow.
	E.g., \expword{\textipa{[p], [k], [b], [m], [n]}}.
	
	\item \optword{fricative}: mouth, pharynx or glottis tightening without there complete closure.
	E.g., \expword{\textipa{[f], [v], [s]}}.
	
	\item \optword{Affricate}: airflow is blocked then released.
	E.g., \expword{\textipa{[\t{\textteshlig}]}}.
	
	\item \optword{lateral}: airflow through a lateral channel (sometimes bilateral).
	E.g., \expword{\textipa{[l]}}.
	
	\item \optword{nasal}: lowering the soft palate.
	E.g., \expword{\textipa{[m], [n]}}.
	
	\item \optword{click}: with the tongue or lips without using the lungs.
	E.g., \expword{\textipa{[!]} (tongue clicking)}.
	
\end{itemize}

\end{frame}

\begin{frame}
\frametitle{Introduction to NLP: Levels}
\framesubtitle{Phonetics, phonology, orthography: Phonetics (IPA)}

\begin{center}
	\vgraphpage{IPA2020_.pdf}
\end{center}

\end{frame}

\begin{frame}
\frametitle{Introduction to NLP: Levels}
\framesubtitle{Phonetics, phonology, orthography: Phonology}


\begin{itemize}
	\item Study of the sounds or phonemes of a given language.
	\item Interested in sounds as elements of a system.
\end{itemize}

\begin{exampleblock}{E.g.: Phoneme \textit{/r/}}
	\begin{itemize}
		\item In French, \textit{r} can be pronounced (in phonetics) as: Trill (rolled R) \expword{\textipa{[r]}}, Voiced uvular trill \expword{\textipa{[\;R]}}, or Voiced uvular fricative (Parisian) \expword{\textipa{[K]}}.
		\item It is always transcribed in the same manner, e.g. \expword{rat /rat/}.
		\item In Arabic, consonants \expword{\RL{r} \textipa{[r]}} and \expword{\RL{.g} \textipa{[G][K]}} have two different phonemes: \expword{\textipa{/r/}} and \expword{\textipa{/G/ /K/}} respectively. 
		E.g., \expword{\RL{.gryb} \textipa{/G\ae ri:b/} (stranger)}.
	\end{itemize}
\end{exampleblock}

\end{frame}

\begin{frame}
\frametitle{Introduction to NLP: Levels}
\framesubtitle{Phonetics, phonology, orthography: Orthography}


\begin{itemize}
	\item Study of the types and form of lemmas/morphemes.
	\item Writing system (according to the \keyword{grapheme}): 
	\begin{itemize}
		\item \optword{logographic}: logograms, each is a unique grapheme denoting a lemma (word).
		E.g., \expword{
			Kanji (Japanese) : 
			\begin{CJK}{UTF8}{min}日, 本, 語\end{CJK}
		}.
		\item \optword{syllabic}: symbols, each represents a syllable (vocalized sound).
		E.g., \expword{
			Hiragana (Japanese) : 
			\begin{CJK}{UTF8}{min}る, た, め, の\end{CJK} ; 
			Katakana (Japanese) : 
			\begin{CJK}{UTF8}{min}セ, ク\end{CJK} ; 
		}.
		\item \optword{alphabetic}: letters; each represents a phoneme.
		E.g., \expword{Latin: A, B, C, etc. }
		\item \optword{abjad}: like the alphabetic writing system, but there are only consonants. 
		E.g. \expword{Arabic: \RL{b}, \RL{t}, \RL{h}}.
	\end{itemize}
	\item Punctuation.
	\item Writing rules.
\end{itemize}

\end{frame}

\begin{frame}
\frametitle{Introduction to NLP: Levels}
\framesubtitle{Phonetics, phonology, orthography: Some humor}

\begin{center}
	\vgraphpage{humor/humor-phonetics.jpg}
\end{center}

\end{frame}

\subsection{Morphology and syntax}

\begin{frame}
\frametitle{Introduction to NLP: Levels}
\framesubtitle{Morphology and syntax: Morphology}

\begin{itemize}
	\item Study of word formation, including how new words are invented in languages.
	\item Study of how words' forms vary according to their uses in sentences.
	\item \keyword{Morpheme}: the smallest unit of language that has its own meaning. 
	E.g., \expword{proper nouns, suffixes, etc.}
	\item \keyword{Lexeme}: a set of all grammatical forms that have the same meaning.
	\item \keyword{Lemma}: a chosen word to represent the lexeme.
	\item Parts of speech: open class (\optword{adjective}, \optword{noun}, \optword{verb}) et closed class (\optword{adverb}, \optword{article}, \optword{conjunction}, \optword{interjection},  \optword{preposition}, \optword{pronoun}).
\end{itemize}

\end{frame}


\begin{frame}
\frametitle{Introduction to NLP: Levels}
\framesubtitle{Morphology and syntax: Morphology}


\begin{itemize}
	\item \optword{Isolating/analytical languages}: each word consists of one and only one morpheme. 
	Morphological changes are few or even absent.
	Among these languages: \expword{Mandarin, Vietnamese, Tai, Khmer, etc.}. 
	E.g., \expword{\begin{CJK}{UTF8}{min}四个男孩\end{CJK} /sì ge nánhái/ ``four boys" (lit. ``four [entity of] masculine child")}.
	\item \optword{Inflectional/synthetic languages}: words are formed from a \keyword{root} plus additional morphemes.
	\begin{itemize}
		\item \optword{Agglutinating languages}: morphemes are clearly phonetically distinguishable from each other.
		Among these languages: \expword{Finnish, Turkish, Japanese, Berber languages, etc.}. 
		E.g., \expword{\begin{CJK}{UTF8}{min}行く, 行きます\end{CJK}}.
		\item \optword{Fusional languages}: it is not always easy to distinguish morphemes from root, or from each other.
		Among these languages: \expword{Arabic, English, French, etc.}.
		E.g., \expword{\RL{kitAb, kutub, 'wa'`.taynAkumuwh}, foot, feet}
	\end{itemize}
\end{itemize}

\end{frame}

\begin{frame}
\frametitle{Introduction to NLP: Levels}
\framesubtitle{Morphology and syntax: Inflectional morphology}


\begin{itemize}
	\item Forming words without changing their categories or creating new lexemes.
	\item Modifying lexemes form so that they adapt to different grammatical contexts.
	\item \optword{Flexion}: \optword{declension} or \optword{conjugation}.
	\item \optword{Affixation}:
	\begin{itemize}
		\item \optword{Prefix}: \expword{\RL{_dhb}} (past), \expword{\RL{y_dhb}} (present).
		\item \optword{Infix}
		\item \optword{Suffix} \expword{étudiant} (singular-masculine), \expword{étudiantes} (plural-feminine).
	\end{itemize}
	\item \optword{Reduplication}: \expword{super-duper, bye-bye, \begin{CJK}{UTF8}{min}きらきら\end{CJK}}.
	
\end{itemize}
\end{frame}

\begin{frame}
\frametitle{Introduction to NLP: Levels}
\framesubtitle{Morphology and syntax: Inflectional morphology (grammatical features)}

\begin{itemize}
	\item \optword{Number}: singular, dual, trial, paucal, plural, etc. 
	\item \optword{Person}: first, second, third, etc.
	\item \optword{Gender}: masculine, feminine, neuter, common.
	\item \optword{Case}: nominative, absolutive, accusative, ergative, etc.
	\item \optword{Tense}: past, present, future.
	\item \optword{Aspect}: accomplished/unaccomplished, progressive, perfective/imperfective, iterative, prospective.
	\item \optword{Voice}: active, medium, passive, etc.
	\item \optword{Polarity}: affirmative, negative.
	\item \optword{Politeness}: informal, formal, etc.
\end{itemize}

\end{frame}

\begin{frame}
\frametitle{Introduction to NLP: Levels}
\framesubtitle{Morphology and syntax: Derivational morphology}


\begin{itemize}
	\item Forming words by changing part of speech (\expword{play, player}) or by creating new lexemes (\expword{do, undo}).
	
	\item \optword{Affixation}: using prefixes, infixes and/or suffixes. 
	E.g., \expword{happy (ADJ), unhappy (ADJ), unhappiness (N)}; \expword{\RL{jhd} (V), \RL{ijthd} (V)}.
	
	\item \optword{Compounding}: by merging words into one.
	E.g., \expword{wind (N) + mill (N) = windmill (N)}.
	
	\item \optword{Blending}: by merging fragments of words into one.
	E.g., \expword{porter (V) + manteau (N) = portemanteau (N)}; \expword{breakfast (N) + lunch (N) = brunch (N)}.
	
	\item \optword{Conversion}: by changing a word's part of speech without any modification. 
	E.g., \expword{orange (fruit, N), orange (color, ADJ)}; \expword{visit-er (V), visite (N)}; \expword{fish (N), to fish (V)}.
	
	\item \optword{Truncation}:
	E.g., \expword{bibliographie, biblio}, \expword{information, info}.
	
	\item \optword{Reduplication}: E.g., \expword{\RL{kr} (V), \RL{krkr} (V)}.
	
\end{itemize}

\end{frame}

\begin{frame}
\frametitle{Introduction to NLP: Levels}
\framesubtitle{Morphology and syntax: Syntax}

\begin{itemize}
	\item Sentence structure: how words are combined to form sentences.
	\item Parts of speech (verb, noun, adjective, etc.) of words in a sentence (PoS tagging).
	\item Word order according to \keyword{Subject (S)}, \keyword{Verb (V)} and \keyword{Object (O)}.
	\item Grammatical functions: \expword{Subject, Direct object, Indirect object, etc.}
\end{itemize}

\end{frame}

\begin{frame}
\frametitle{Introduction to NLP: Levels}
\framesubtitle{Morphology and syntax: Syntax (Word order)}

\begin{table}
	\begin{tblr}{
			colspec = {p{.1\textwidth}p{.15\textwidth}p{.65\textwidth}},
			row{odd} = {lightblue},
			row{even} = {lightyellow},
			row{1} = {darkblue},
		} 
		\textcolor{white}{Order} & \textcolor{white}{Proportion} & \textcolor{white}{Examples} \\
		SOV & 44.78\% & Japanese, Latin, Tamil, Basque, Urdu, Ancient Greek, Bengali, Hindi, Sanskrit, Persian, Korean \\
		SVO & 41.79\% & French, Mandarin, Russian, English, Hausa, Italian, Malay (language), Spanish, Thai\\
		VSO & 9.20\% & Irish, Arabic, Biblical Hebrew, Filipino, Tuareg languages, Welsh \\
		VOS & 2.99\% & Malagasy, Baure, Car (language) \\
		OVS & 1.24 \% & Apalai, Hixkaryana, Klingon (language) \\
	\end{tblr}
	\caption{Words Order proportions based on 402 languages \cite{1988-blake}}
\end{table}

\end{frame}

\begin{frame}
\frametitle{Introduction to NLP: Levels}
\framesubtitle{Morphology and syntax: Syntax (Word order, examples)}

\begin{itemize}
	\item \textcolor{blue}{Subject}, \textcolor{red}{Verb}, \textcolor{green}{Object}
\end{itemize}

\begin{columns}
\begin{column}{0.54\textwidth}
	\begin{exampleblock}{Example SOV [Japanese]}
		\small
		\begin{CJK}{UTF8}{min}
			\textcolor{blue}{カリムさん}は\textcolor{green}{日本語}を\textcolor{red}{勉強します}。
		\end{CJK}
	\end{exampleblock}
\end{column}
%
\begin{column}{0.42\textwidth}
	\begin{exampleblock}{Example SVO [English]}
		\textcolor{blue}{Karim} \textcolor{red}{learns} \textcolor{green}{English}.
	\end{exampleblock}
\end{column}%
\end{columns}
\begin{columns}
\begin{column}{0.49\textwidth}
	\begin{exampleblock}{Example VSO [Arabic]}
		\begin{flushright}
			\color{black}.
			\color{green} \RL{al-`rbyT}%
			\color{blue} \RL{krym} 
			\color{red} \RL{yt`llm} 
		\end{flushright}
	\end{exampleblock}
\end{column}
\end{columns}

\end{frame}

\begin{frame}
\frametitle{Introduction to NLP: Levels}
\framesubtitle{Morphology and syntax: Syntax (Constituency grammar)}

\begin{itemize}
	\item A \keyword{sentence} is composed of several \keyword{phrases} which includes words and other phrases.
	\item A phrase contains a \keyword{head} (\keyword{nucleus}) which is the central element.
	\item Depending on its head, a phrase can be: nominal (\keyword{NP}), adjectival (\keyword{AP}), verbal (\keyword{VP}) or prepositional (\keyword{PP}).
	\item The most used formal system to model the constituency structure of a sentence is the \keyword{context-free grammar} (Level 2 in \keyword{Chomsky}'s hierarchy).
\end{itemize}


\end{frame}


\begin{frame}
\frametitle{Introduction to NLP: Levels}
\framesubtitle{Morphology and syntax: Syntax (Constituency grammar: example)}

\begin{columns}
	\begin{column}{0.52\textwidth}
	The grammar that generated this syntactic tree can be written as: 
	
	\begin{enumerate}
		\item S $ \rightarrow $ NP VP
		\item NP $ \rightarrow $ DET NP'
		\item NP $ \rightarrow $ DET N
		\item NP' $ \rightarrow $ ADJ N
		\item VP $ \rightarrow $ V NP
	\end{enumerate}

	The second rule is not written as \\(NP $ \rightarrow $ DET NP)\\ 
	otherwise there can be several determiners to a noun.
	
	\end{column}
	%
	\begin{column}{0.42\textwidth}
		\hgraphpage{const_gram.pdf}
	\end{column}%
\end{columns}


\end{frame}

\begin{frame}
\frametitle{Introduction to NLP: Levels}
\framesubtitle{Morphology and syntax: Syntax (Dependency grammar)}

\begin{itemize}
	\item Syntactic structure is described in terms of words and not phrases.
	\item \optword{Dependency grammar}: a set of binary relations between the words of a sentence.
	\item The relations can be: a nominal subject (\optword{nsubj}), an object (\optword{obj}), an adjective modifier (\optword{amod}), determiner (\optword{det}), etc.
\end{itemize}

\begin{figure}
	\centering
	\hgraphpage[0.6\textwidth]{dep_gram_.pdf}
	\caption{Dependencies example, generated using \url{https://corenlp.run/}}
\end{figure}

\end{frame}

\begin{frame}
\frametitle{Introduction to NLP: Levels}
\framesubtitle{Morphology and syntax: Some humor}

\begin{center}
	\vgraphpage{humor/humor-syntax.jpg}
\end{center}

\end{frame}

\subsection{Semantics}

\begin{frame}
\frametitle{Introduction to NLP: Levels}
\framesubtitle{Semantics}

\begin{itemize}
	\item Study of meaning.
	\begin{itemize}
		\item \optword{Lexical semantics}: words meaning.
		\item \optword{Propositional semantics}: sentences meaning.
	\end{itemize}
\end{itemize}

\begin{exampleblock}{Example of different meanings}
	The word \expword{sound} has many meanings (from \url{http://wordnetweb.princeton.edu/perl/webwn})
	\begin{itemize}
		\item The cat is chasing some \expword{chickens}. ``\textit{a domestic fowl bred for flesh or eggs.}"
		\item I ate some \expword{chicken}. ``\textit{the flesh of a chicken used for food.}"
		\item Don't be a \expword{chicken}! ``\textit{A person who lacks confidence, is irresolute and wishy-washy.}"
	\end{itemize}
	
\end{exampleblock}

\end{frame}

\begin{frame}
\frametitle{Introduction to NLP: Levels}
\framesubtitle{Semantics: Words meaning}

\begin{itemize}
	\item In logographic systems, a \keyword{grapheme} represents one or more meanings (usually with multiple pronunciations).
	
	\begin{CJK}{UTF8}{min}
		\expword{川} (river), \expword{山} (mountain), 
		
		\expword{音} (sound, noise) + \expword{楽} (music, comfort, ease) = \expword{音楽} (music)
	\end{CJK}

	\item A word can have several meanings (\keyword{Polysemy}).
	\item The meaning of a word is built...
	\begin{itemize}
		\item of a set of ``semantic primitives".
		\item using relations between morphemes (a network of meanings).
	\end{itemize}
\end{itemize}

\end{frame}


\begin{frame}
\frametitle{Introduction to NLP: Levels}
\framesubtitle{Semantics: Words meaning (Semic analysis)}

\begin{itemize}
	\item Proposed by Eric Buyssens in the 1930s and developed by Bernard Pottier in the 1960s.
	\item \keyword{Seme} a minimal unit of meaning (minimal semantic feature).
	\item \keyword{Sememe} a seme cluster corresponding to a lexical unit.
	\item Words having one or more positive semes in common belong to the same semantic field.
\end{itemize}

\begin{exampleblock}{Example of semic analysis}
	\centering
	\begin{tabular}{|l|l|l|l|}
		\hline
		Word/Seme & animate & domestic & feline \\
		\hline
		Cat & + & + & + \\
		\hline
		Lion & + & - & + \\
		\hline
		Dog & + & + & - \\
		\hline
	\end{tabular}
\end{exampleblock}

\end{frame}

\begin{frame}
\frametitle{Introduction to NLP: Levels}
\framesubtitle{Semantics: Words meaning (Semantic relations)}

\begin{itemize}
	\item \optword{Synonymy}: if a word is substituted for another in a sentence without changing the meaning, then these two words are synonyms.
	\item \optword{Antonymy}: the opposite meaning. The two words must express two values of the same property. E.g., \expword{big} and \expword{small} express the property \expword{size}.
	
	\item \optword{Hyponymy}: a word with a more specific meaning than another.
	E.g., \expword{cat} is the hyponym of \expword{feline} hyponym of \expword{animal}.
	 
	\item \optword{Hypernymy}: a word with a more generic meaning.
	E.g., \expword{feline} is the hyperonym/hypernym of \expword{cat}.
	
	\item \optword{Meronymy}: a word that denotes a part of another.
	E.g., \expword{wheel} is a meronym of \expword{car}.
\end{itemize}

\end{frame}

\begin{frame}
\frametitle{Introduction to NLP: Levels}
\framesubtitle{Semantics: Propositions meaning}

\begin{minipage}{0.5\textwidth}
\begin{itemize}
	\item First-order logic.
	\item \keyword{predicate} is a property of entities. 
	E.g. \expword{Possess, Person}.	
	\item \keyword{constant} a specific world entity.
	E.g. \expword{Karim, Algeria, ESI}.
	\item \keyword{variable} denotes an anonymous entity. 
	E.g. \expword{x, y}.
	\item \keyword{function} results in a specific entity without the need to create a new constant.
	E.g. \expword{locationOf : locationOf(ESI) vs. locationOfESI}.
\end{itemize}
\end{minipage}
\begin{minipage}{0.48\textwidth}
	\begin{figure}
		{\tiny\bfseries
		\begin{tabular}{rcl}
			\hline\hline
			Formula & \textrightarrow & AtomicFormula \\
			& \textbar        & Formula Connective Formula \\
			& \textbar        & Quantifier Variable, ... Formula \\
			& \textbar        & $\textlnot$ Formula \\
			& \textbar        & (Formula) \\
			AtomicFormula & \textrightarrow & Predicate (Term, ...) \\
			Term    & \textrightarrow & Function(Term, ...) \\
			& \textbar        & Constant \\
			& \textbar        & Variable \\
			Connective & \textrightarrow & $\wedge$ \textbar $\vee$ \textbar $\Rightarrow$ \\
			Quantifier & \textrightarrow & $\forall$ \textbar $\exists$ \\
			Constant & \textrightarrow & ESI \textbar Karim \textbar Algérie ...\\
			Variable & \textrightarrow & x \textbar y \textbar ... \\
			Predicate & \textrightarrow & Ecole \textbar EnseignatA \textbar Utiliser \textbar ... \\
			Function & \textrightarrow & EmplacementDe \textbar ... \\
			\hline\hline
		\end{tabular}
	}
		\caption{Syntax adapted from \cite{2002-russell-norvig} describing the representation of first-order logic \cite{2019-jurafsky-martin}}
	\end{figure}
\end{minipage}

\end{frame}

\begin{frame}
\frametitle{Introduction to NLP: Levels}
\framesubtitle{Semantics: Propositions meaning (Example)}

\begin{exampleblock}{Exemple de la représentation sémantique d'une phrase}
	\begin{itemize}
		%	\item I have a car \cite{2019-jurafsky-martin}
		%	\[%
		%	\exists\, e, y\, Having(e) \wedge Haver(e, Speaker) \wedge HadThing(e, y) \wedge Car(y)
		%	\]
		\item Each of some students owns two computers.
		\item Two computers are owned by each of some students.
		\item $S = Student, \; C = Computer, \; O = Own$
		\item $\exists x (S(x) \wedge \forall y ( C(y) \Rightarrow O(x, y)) )$ \textcolor{red}{\XBox}
		\item $\exists x (S(x) \wedge \exists y, z (( (C(y) \wedge O(x, y) ) \wedge (C(z) \wedge O(x, z) ) ))$ \textcolor{red}{\XBox}
		\item $\exists x (S(x) \wedge \exists y, z (\neg y = z \wedge ( (C(y) \wedge P(x, y) ) \wedge (C(z) \wedge O(x, z) ) ))$ \textcolor{green}{\CheckedBox}
		\item $\exists x (S(x) \wedge \exists y ((C(y) \wedge O(x, twoOf(y)) ))$ \textcolor{green}{\CheckedBox}
	\end{itemize}
\end{exampleblock}

\end{frame}

\begin{frame}
\frametitle{Introduction to NLP: Levels}
\framesubtitle{Semantics: Propositions meaning (Inference)}

\begin{exampleblock}{Example of inference}
	\centering
	\begin{tabular}{lll}
		Each human is mortal.  & & Socrates is a human. \\
		$\forall x (Human(x) \Rightarrow Mortal(x))$ && $Human(Socrates)$ \\
		\hline
		\multicolumn{3}{c}{$Mortal(Socrates)$}\\
	\end{tabular}
	
\end{exampleblock}

\end{frame}

\begin{frame}
\frametitle{Introduction to NLP: Levels}
\framesubtitle{Semantics: Some humor}

\begin{center}
	\vgraphpage{humor/humor-semantics.jpg}
\end{center}

\end{frame}

\subsection{Pragmatics and discourse}

\begin{frame}
\frametitle{Introduction to NLP: Levels}
\framesubtitle{Pragmatics and discourse: Pragmatics (Motivation)}

\begin{minipage}{0.7\textwidth}
\begin{itemize}
	\item Imagine: we are in an anniversary;
	\item A person shouted: ``Candles!" ;
	\item We can immediately understand that there are no candles on the pie.
	\item \optword{Question}: how can we infer that?
	\item \optword{Answer}: from \keyword{context}.
\end{itemize}
\end{minipage}
\begin{minipage}{0.28\textwidth}
	\hgraphpage{pragmatics.pdf}
\end{minipage}

\vfill
\begin{tabular}{|p{.4\textwidth}|p{.5\textwidth}|}
	\hline
	\bfseries Semantics & \bfseries Pragmatics \\
	\hline 
	\multicolumn{2}{|l|}{A person \textbf{A} arrived late for an appointment.} \\
	\multicolumn{2}{|l|}{\textbf{B}: ``What time do you think it is?"} \\
	\hline 
	\textbf{A} announces the time &  \textbf{A} explains the reason for the delay \\
	\hline
\end{tabular}

\vfill 

\optword{Humor}: \expword{``I'm sorry" and ``My bad" have the same meaning ... Unless you are in a funeral}

\end{frame}

\begin{frame}
\frametitle{Introduction to NLP: Levels}
\framesubtitle{Pragmatics and discourse: Pragmatics (Reasoning)}

\begin{itemize}
	\item \optword{Implicature}: refers to what the speaker means implicitly.
	According to \cite{1979-Grice}, the interlocutors must respect certain conversational norms (maxims): quantity, quality, relevance and manner.
	
	E.g. \expword{Karim, who is lazy, has stopped writing. \textrightarrow\ 
		Karim was writing. Karim is lazy. Karim no longer writes. There is an entity called Karim.}
	
	\item \optword{Presupposition}: refers to the assumptions made by the interlocutors during a communication.
	
	\item \optword{Speech act}: refers to the linguistic interaction of the speaker in order to act on his environment. A list of speech acts is provided by \cite{1962-austin}: declaration, order, question, prohibition, greeting, invitation, congratulation, apology.
\end{itemize}

\end{frame}

\begin{frame}
\frametitle{Introduction to NLP: Levels}
\framesubtitle{Pragmatics and discourse: Discourse (Coreference)}

\begin{exampleblock}{Example of coreference}
	\begin{itemize}
		\item \keyword{The cat} doesn't fit in the box because \keyword{it} is too big.
		\item The cat doesn't fit in \keyword{the box} because \keyword{it} is too small.
	\end{itemize}
\end{exampleblock}

\begin{itemize}
	\item \optword{Pronouns}: \expword{the cat} chased a mouse and \expword{it} is playing with it.
	\item \optword{Nominal phrases}: \expword{Apple} is a computer manufacturer. \expword{The firm} is known worldwide.
	\item \optword{Zero Anaphora}: in languages like Japanese, sometimes the reference is omitted.
	\item \optword{Nouns}: \expword{International Business Machine} is an American company. Like Apple, \expword{IBM} manufactures computers.
\end{itemize}

\end{frame}

\begin{frame}
\frametitle{Introduction to NLP: Levels}
\framesubtitle{Pragmatics and discourse: Discourse (Coherence 1)}

\begin{itemize}
	\item Logical links between units (sentences and propositions) of the text.
	\item \keyword{Rhetorical Structure Theory (RST)}: the most widely used discourse model.
	\item Two types of clauses: \optword{nuclear} and \optword{satellite}
	\begin{itemize}
		\item \expword{The student was absent yesterday. He was sick.}
		\item First sentence is a nucleus since it describes the main event.
		\item Second sentence is a satellite since it depends on the first.
	\end{itemize}
	\item There are several coherence relations: reason, elaboration, evidence, attribution, narration, etc.
\end{itemize}

\end{frame}

\begin{frame}
\frametitle{Introduction to NLP: Levels}
\framesubtitle{Pragmatics and discourse: Discourse (Coherence 2)}

\begin{exampleblock}{Example of "reason"}
	[\textsubscript{N} The student was absent yesterday.] [\textsubscript{S} He was sick.]
\end{exampleblock}

\begin{exampleblock}{Example of "elaboration"}
	[\textsubscript{N} The exam is easy.] [\textsubscript{S} It takes only one hour.]
\end{exampleblock}

\begin{exampleblock}{Example of "evidence"}
	[\textsubscript{N} Kevin must be here.] [\textsubscript{S} His car is parked outside.]
\end{exampleblock}

\end{frame}

\begin{frame}
\frametitle{Introduction to NLP: Levels}
\framesubtitle{Pragmatics and discourse: Some humor}

\begin{center}
	\vgraphpage{humor/humor-pragmatics.jpg}
\end{center}

\end{frame}

%===================================================================================
\section{Applications of NLP}
%===================================================================================

%\begin{frame}
%\frametitle{Introduction au TALN}
%\framesubtitle{Applications du TALN}
%
%\end{frame}

\subsection{Tasks}

\begin{frame}
\frametitle{Introduction to NLP: Applications}
\framesubtitle{Tasks: Morphosyntax}

\begin{itemize}
	\item \optword{Sentence Boundary Disambiguation and Words Tokenization}: breaking text into smaller units for easier processing.
	\item \optword{Stemming and lemmatization}: reducing inflectional (sometimes also derivational) forms of a word to a common base form.
	\item \optword{Part of Speech tagging}: finding parts of speech of a sentence's words.
	\item \optword{Parsing}: finding the syntactic structure of a sentence (how it was formed).
	\item \optword{Terminology extraction}: searching for existing terminologies in a text.
\end{itemize}

\end{frame}

\begin{frame}
\frametitle{Introduction to NLP: Applications}
\framesubtitle{Tasks: Semantics}

\begin{itemize}
	\item \optword{Word-sense disambiguation}: identifying a a word's meaning.
	\item \optword{Semantic role labeling}: assigning semantic roles (agent, theme, etc.) to words or phrases in a sentence.
	\item \optword{Named Entity Recognition}: extracting names of people, names of organizations or companies, names of places, quantities, distances, values, dates, etc.
	\item \optword{Semantic analysis}: finding the semantic representation of a text.
	\item \optword{Paraphrase}: formulating a text differently.
	\item \optword{Entailment}: checking if a segment of the text is true implies that another is also true.
	\item \optword{Automatic text generation}.
\end{itemize}

\end{frame}

\begin{frame}
\frametitle{Introduction to NLP: Applications}
\framesubtitle{Tasks: Discourse}

\begin{itemize}
	\item \optword{Coreference resolution}: linking the different references in a text.
	\item \optword{Coherence}: looking for relations between sentences.
	\item \optword{Ellipsis resolution}: finding omitted elements of text.
	Example of an ellipsis: \expword{``Peter eats cherries, Paul strawberries"}.
\end{itemize}

\end{frame}

\subsection{Systems}

\begin{frame}
\frametitle{Introduction to NLP: Applications}
\framesubtitle{Systems: Machine translation}

\begin{itemize}
	\item Translation of a text from a language to another.
\end{itemize}

\begin{figure}
	\hgraphpage{exp.MT.png}
	\caption{Example of Google translate.}
\end{figure}

\end{frame}

\begin{frame}
\frametitle{Introduction to NLP: Applications}
\framesubtitle{Systems: Automatic text summarization}

\begin{itemize}
	\item Generating a compact version of a text.
\end{itemize}

\begin{figure}
	\hgraphpage{exp.ATS.png}
	\caption{Example of esummarizer.}
\end{figure}

\end{frame}

\begin{frame}
\frametitle{Introduction to NLP: Applications}
\framesubtitle{Systems: Question/Answering}

\begin{itemize}
	\item Affording answers to questions formulated in natural language.
\end{itemize}

\begin{figure}
	\hgraphpage{exp.QR.png}
	\caption{Example of START.}
\end{figure}

\end{frame}

\begin{frame}
\frametitle{Introduction to NLP: Applications}
\framesubtitle{Systems: Conversational agents (chatbots)}

\begin{itemize}
	\item Chatting with a user.
\end{itemize}

\begin{figure}
	\centering
	\hgraphpage[.5\textwidth]{exp.chatbot.png}
	\caption{Example of Cleverbot.}
\end{figure}

\end{frame}

\begin{frame}
\frametitle{Introduction to NLP: Applications}
\framesubtitle{Systems: Information extraction}

\begin{itemize}
	\item \optword{Text mining}: extraction of knowledge from a text.
	\item \optword{Sentiment analysis}: finding existing sentiments in a text.
	E.g., \expword{checking whether users are satisfied with a product or not}.
	\item \optword{Automatic document recommendation}: presenting items that are likely to interest a user.
	E.g., \expword{recommend books}.
\end{itemize}

\end{frame}

\begin{frame}
\frametitle{Introduction to NLP: Applications}
\framesubtitle{Systems: Memes generation}

\begin{figure}
	\centering
	\vgraphpage[0.7\textheight]{exp.meme.png}
	\caption{Architecture to generate memes \cite{sadasivam2020memebot}.}
\end{figure}

\end{frame}


\subsection{Business}

\begin{frame}
\frametitle{Introduction to NLP: Applications}
\framesubtitle{Business: Commerce}

\begin{itemize}
	\item \optword{Publicity}: identifying new audiences potentially interested in certain products.
	\item \optword{Customer Service}: using chatbots to answer potential customer questions.
	Also, using sentiment analysis to get an idea about customers opinions on the company's products.
	\item \optword{Market Intelligence}: monitoring blogs, websites and social networks to analyze market trends and get an idea of the competition.
	\item \optword{Recruitment}: filtering CVs to find candidates faster and without bias.
\end{itemize}

\end{frame}

\begin{frame}
\frametitle{Introduction to NLP: Applications}
\framesubtitle{Business: E-Governance}

\begin{itemize}
	\item \optword{Government/Citizen Communication}: Illiterate citizens can share their opinions using audio/video, which can be translated into text. Similarly, a message to citizens can be transformed into a voice message.
	\item \optword{Opinion mining}: extraction of comments, complaints and criticisms on a particular policy, thus determining the general consensus on this subject.
\end{itemize}

\end{frame}

\begin{frame}
\frametitle{Introduction to NLP: Applications}
\framesubtitle{Business: Health}

\begin{itemize}
	\item Structuring medical documents to facilitate their use.
	\item Researching, analyzing and interpreting massive amounts of patient datasets.
	\item Predicting diseases based on symptoms and medical records.
	\item Generating reports.
	\item Using virtual assistants to monitor and assist patients.
\end{itemize}

\end{frame}

\begin{frame}
\frametitle{Introduction to NLP: Applications}
\framesubtitle{Business: Education}

\begin{itemize}
	\item Language assessment: reading, writing and speaking.
	\item Spell checking.
	\item Automatic assessment of students work, such as essays and answers.
	\item E-learning by integrating chatbots with games, which promotes an active learning environment.
\end{itemize}

\end{frame}

%===================================================================================
\section{NLP Challenges}
%===================================================================================


\subsection{Resources}

\begin{frame}
\frametitle{Introduction to NLP: Challenges}
\framesubtitle{Resources}

\begin{itemize}
	\item Lack of tools and datasets for less spoken languages.
	\item Manual annotation of training and test corpora.
	\item Large document processing: machine learning is limited when it comes to representing long contexts.
\end{itemize}
\end{frame}

\subsection{Language understanding}

\begin{frame}
\frametitle{Introduction to NLP: Challenges}
\framesubtitle{Language understanding: Ambiguity and Coreference}

\begin{itemize}
	\item \keyword{polysemy}: a word with multiple meanings.
	E.g. \expword{Indian: from India, native to the Americas}.
	\item \keyword{enantiosemy}: a polysemous term with two antonymous meanings.
	E.g. \expword{Regret, denotes complaint or nostalgia: ``I regret my childhood".}
	\item \keyword{Trope} (figurative language): metaphors (\expword{``It's raining cats and dogs"}), metonomy (\expword{``The room applauded" [the people in the room]}), irony (``What a beautiful day!" [to mean that it's raining cats and dogs.]).
	\item Co-references in long texts (ambiguity in reference choice).
	Ex. \expword{\underline{Table data} is dumped into \underline{a delimited text file}, which is sent to \underline{the remote site} where \underline{it} is loaded into the destination database. }
\end{itemize}

\end{frame}

\begin{frame}
\frametitle{Introduction to NLP: Challenges}
\framesubtitle{Language understanding: Human aspect}

\begin{itemize}
	\item \optword{Personality}: how to model a personality? how does personality affect speech?
	\item \optword{Variations and non-standard communication}: users do not respect writing standards of a language. E.g., \expword{chat language, Arabizi, Frenglish, etc.}
	\item \optword{Intention}: there are sentences that do not mean what we understand directly; they mean something else.
	\item \optword{Emotions}: sentences can change meaning depending on the accompanied emotions.
\end{itemize}

\end{frame}


\subsection{Evaluation}

\begin{frame}
\frametitle{Introduction to NLP: Challenges}
\framesubtitle{Evaluation}

\begin{itemize}
	\item \optword{Manual evaluation}
	\begin{itemize}
		\item takes time (high time cost).
		\item sometimes requires experts (high cost in terms of expenses).
		\item problem of evaluator subjectivity and inter-evaluator agreement.
	\end{itemize}
	\item \optword{Automatic evaluation}
	\begin{itemize}
		\item several aspects are difficult to assess automatically.
	\end{itemize}
\end{itemize}

\end{frame}

\subsection{Ethics}

\begin{frame}
\frametitle{Introduction to NLP: Challenges}
\framesubtitle{Ethics}

\begin{itemize}
	\item \optword{Privacy}: collected data may contain private information of individuals. Companies can store information about their users.
	\item \optword{Bias and Discrimination}: Corpora used to train a system may cause bias.
	\item \optword{Usage}: espionage and social engineering, loss of jobs due to automation, etc.
\end{itemize}

\end{frame}

\begin{frame}
\frametitle{Introduction to NLP: Challenges}
\framesubtitle{Ethics: Some Humor}

\begin{center}
	\vgraphpage{humor/humor-ethics.jpg}
\end{center}

\end{frame}

\insertbibliography{NLP01}{*}

\end{document}

